% ==============================================================================
% 05_diskussion.tex – Kapitel 5
% ==============================================================================
\chapter{Diskussion}
\label{ch:diskussion}

\section{Antworten auf die Forschungsfragen}
\label{sec:antworten}

\paragraph{FF1: Wie stark hilft eine kurze Personalisierung?}
Das allgemeine Basismodell ist auf den eigenen Aufnahmen nicht besser als die
Konstant-Baseline (\SI{30,57}{\percent} gegenüber \SI{30,32}{\percent}). Bei
einem zeitlichen Split, der der echten Nutzung entspricht, verbessert die
Personalisierung das Ergebnis. Der Fehler sinkt von \SI{31,84}{\percent} auf
\SI{27,93}{\percent}. Das ist eine Verbesserung von etwa \SI{12}{\percent}
(Wilcoxon $p=0,002$). In 8 von 18 Sitzungen gibt es jedoch keine Verbesserung
(\cref{sec:design_performance}). Bei den guten Sitzungen 21 und 22 ist die
Verbesserung mit \SI{46}{\percent} weiterhin deutlich. Auch bei einer zweiten
Aufnahme bleibt der Gewinn mit \SI{47}{\percent} fast gleich
(\cref{sec:loading}).

Bei einem zufälligen Split sind die Ergebnisse noch besser. Der Fehler sinkt auf
\SI{19,79}{\percent}, also um etwa ein Drittel. Die beste Sitzung erreicht sogar
einen Fehler von nur \SI{10,5}{\percent} (\cref{sec:sessions_results}). Dieser
Split zeigt aber ein zu positives Ergebnis und entspricht nicht der echten
Nutzung.

Bei Person~B brachte die Personalisierung mit dem vollständigen Kamerabild in
beiden Sitzungen keine Verbesserung. Das gilt auch beim zufälligen Split. Mit
dem Augenausschnitt aus der Web-Demo sank der Fehler dagegen von
\SI{26,9}{\percent} auf \SI{9,7}{\percent} (\cref{sec:person_b}).

\emph{Antwort:} Eine kurze Personalisierung kann den Fehler deutlich
reduzieren. Das funktioniert aber nicht bei jeder Aufnahme gleich gut. Gute
Aufnahmebedingungen und ein passender Bildausschnitt der Augen sind dafür
wichtig. Ein zufälliger Split zeigt außerdem bessere Ergebnisse, als man bei der
echten Nutzung erwarten kann.

\paragraph{FF2: Welche Trainingseinstellung eignet sich?}
Auf GazeCapture zeigen die meisten Trainingseinstellungen keinen deutlichen
Unterschied. Die Ergebnisse aller Modelle liegen nahe an der Konstant-Baseline
(\cref{sec:training}). Die Tests zeigen aber, dass Sigmoid als Activation die
besten Ergebnisse liefert. Außerdem ist L1-Loss besser als der Loss mit
Positionsgewichtung. Es reicht aus, nur den Regression Head zu trainieren. Das
liefert ähnliche Ergebnisse wie das Training weiterer Schichten, ist aber
stabiler und deutlich schneller.

\emph{Antwort:} Für das Training eignet sich folgende Einstellung:
\begin{itemize}[nosep]
  \item eingefrorener Backbone (wird nicht trainiert),
  \item Regression Head mit Sigmoid,
  \item L1-Loss,
  \item Learning Rate $10^{-4}$,
  \item Batch Size 32.
\end{itemize}
Diese Einstellung wurde vor allem wegen ihrer Stabilität und des geringen
Trainingsaufwands ausgewählt. Die Learning Rate $10^{-4}$ wurde zusätzlich
gewählt, weil sie im ersten Test den kleinsten Fehler erreicht hat.

\paragraph{FF3: Welche Rolle spielen Datenaufteilung, Kalibrierung und Glättung?}
Die Aufteilung der Daten hat einen großen Einfluss auf die Ergebnisse. Ein
zufälliger Split liefert zu gute Ergebnisse, weil fast gleiche Bilder aus einer
Aufnahme sowohl im Training als auch im Test vorkommen. Die gleichmäßige
Bewegung des Ziels über den Bildschirm (Gen~1) deckt den Bildschirm am besten
ab und erreicht den kleinsten Fehler; der Unterschied ist allerdings nicht
signifikant (\cref{sec:design_performance}). Bei schnellen Sprüngen entstehen
Bilder mit falschen Labels. Wenn diese WAIT-Bilder entfernt werden, verbessert
sich das Ergebnis leicht um \num{0,84} Prozentpunkte. Diese Verbesserung ist
jedoch nicht signifikant. Mehr Daten in den Ecken verbessern die Ergebnisse in
diesen Bereichen nicht. Ein Moving Average über die Vorhersagen macht die
Anzeige deutlich ruhiger. Bei einem Fenster von $k=3$ sinkt die RMS-S2S um
\SI{40}{\percent}. Die Genauigkeit wird dadurch aber nicht besser. Bei längeren
Fenstern wird das Ergebnis sogar schlechter, weil die Vorhersage stärker
verzögert wird (\cref{sec:smoothing}).

\emph{Antwort:} Für die Kalibrierung ist eine kurze und gleichmäßige Bewegung
über den gesamten Bildschirm am besten. Schnelle Sprünge sollten erkannt und
protokolliert werden, damit die Bilder während dieser Bewegungen später
entfernt werden können. Für die Bewertung sollte immer ein zeitlicher Split
verwendet werden, da dieser der echten Nutzung besser entspricht. Ein kurzer
Moving Average (etwa \SI{0,5}{\second}) kann die Anzeige ruhiger machen. Er
verbessert aber nicht die Genauigkeit und ersetzt keine bessere Kalibrierung.

\paragraph{FF4: Welchen Einfluss haben die Aufnahmebedingungen?}
Der größte Einfluss kommt vom Hintergrund. Ein unruhiger Hintergrund erhöht den
Fehler um 6 bis 15 Prozentpunkte. Auch eine Brille führt meistens zu einem
höheren Fehler. In dunklen Räumen hilft die Personalisierung nur wenig. Die
Dauer der Aufnahme und das verwendete Gerät (Webcam) zeigen dagegen keinen
klaren Einfluss (\cref{sec:conditions}).

\emph{Antwort:} Ein ruhiger Hintergrund und gute Lichtverhältnisse sind für gute
Ergebnisse besonders wichtig. Sie haben einen größeren Einfluss als die Art der
Kalibrierung. Diese Ergebnisse basieren jedoch nur auf einzelnen Sitzungen und
müssen deshalb vorsichtig betrachtet werden.

\section{Einordnung in den Stand der Technik}
\label{sec:einordnung}
Einige Ergebnisse passen gut zur Literatur:
\begin{itemize}
  \item \textbf{Wenige Parameter reichen.} Wie bei der Offset-Kalibrierung
        \cite{chen2020offset}, bei SPAZE \cite{linden2019spaze} und bei PIGaze
        \cite{wang2026pigaze} genügt hier ein kleiner, angepasster Head auf
        einem eingefrorenen Backbone. Mehr trainierbare Schichten machen das
        Training nur instabiler. Das passt zu Kumar et
        al.~\cite{kumar2022finetuning}.
  \item \textbf{Personalisierung ist der wichtigste Faktor.} Wie bei
        Valliappan et al.~\cite{valliappan2020smartphone} kommt fast der ganze
        Gewinn aus wenigen Sekunden persönlicher Daten.
  \item \textbf{Das Evaluationsprotokoll entscheidet.} Ähnlich wie beim
        Unterschied zwischen \code{norm\_random} und \code{norm\_subject}
        überschätzen zufällige Splits die Leistung
        (\cref{subsec:norm_subject}).
  \item \textbf{Glättung verbessert die Precision, nicht die Accuracy.} Das
        entspricht der Unterscheidung von Niehorster et
        al.~\cite{niehorster2026quality}: Ein Moving Average verringert die
        Streuung von Bild zu Bild deutlich, der systematische Fehler bleibt
        aber fast gleich (\cref{sec:smoothing}).
  \item \textbf{Falsche Labels haben einen kleinen Effekt.} Auch DMRF
        \cite{jiang2026dmrf} senkt den Fehler nur wenig
        ($3,92^\circ \rightarrow 3,85^\circ$). Hier bringt das Entfernen der
        WAIT-Bilder \num{0,84} Prozentpunkte; dieser Unterschied ist nicht
        signifikant ($p=0,43$).
\end{itemize}
Es gibt aber auch einen wichtigen Unterschied: Das allgemeine Modell dieser
Arbeit ist sehr schwach. Es bleibt bei der Bildschirmmitte, während
Verfahren in der Literatur auch ohne Personalisierung einige Grad Genauigkeit
erreichen. Der wichtigste Grund ist wahrscheinlich die Eingabe: Die meisten
Verfahren nutzen Gesichts- oder Augenausschnitte, diese Arbeit dagegen das
ganze Kamerabild. Außerdem nutzt das Basismodell nur \num{1000} Bilder von 179
Personen. Die absoluten Zahlen sind wegen der unterschiedlichen Metriken
ohnehin nicht direkt mit der Literatur vergleichbar (\cref{sec:vergleich}).

\section{Grenzen der Arbeit}
\label{sec:grenzen}
\begin{itemize}
  \item \textbf{Wenige Personen:} Die 18 ausgewerteten Sitzungen stammen von
        einer Person (dem Autor). Von einer zweiten Person gibt es nur zwei
        Sitzungen und einen Durchlauf der Demo. Dass der Augenausschnitt die
        Ursache des Unterschieds ist, ist ein Hinweis aus diesem einen
        Durchlauf und nicht systematisch geprüft. Die Ergebnisse lassen sich
        daher nicht ohne Weiteres auf andere Personen übertragen.
  \item \textbf{Wenige Sitzungen und vermischte Bedingungen:} 18 Sitzungen
        reichen nicht für sichere Aussagen über einzelne Faktoren. Gestaltung,
        Licht, Hintergrund und Brille ändern sich oft gleichzeitig.
  \item \textbf{Einzelne Läufe:} Pro Sitzung gibt es meist nur einen Lauf
        (Seed 42). Die Streuung zwischen Seeds kann bis zu 3 Prozentpunkte
        betragen.
  \item \textbf{Ein Backbone und ganze Bilder:} Untersucht wurde nur ResNet-18
        mit ganzen Kamerabildern. DINOv2, CLIP, MAE, UniGaze, Adapter und LoRA
        wurden nicht getestet, ebenso wenig Gesichts- oder Augenausschnitte in
        den Experimenten von Kapitel~4.
  \item \textbf{Schwaches Basismodell:} Das Basismodell lernt auf GazeCapture
        kaum etwas und dient in dieser Arbeit vor allem als Startpunkt für die
        Personalisierung. Ob ein zufällig initialisierter Head ähnlich gut
        wäre, wurde nicht getestet. Ein besseres Basismodell könnte die
        Personalisierung verändern.
  \item \textbf{Viele Tests:} Es wurden viele statistische Tests ohne Korrektur
        für Mehrfachvergleiche (zum Beispiel nach Bonferroni) gerechnet.
        Knappe Ergebnisse wie $p=0,032$ (L1 gegenüber dem sensitiven Loss) sind
        daher vorsichtig zu lesen.
  \item \textbf{Glättung nur nachträglich:} Der Moving Average wurde auf
        gespeicherten Vorhersagen des sequenziellen Splits (meist ein Seed)
        und der zwei Videos ausgewertet. Für Fixationen gibt es nur einen
        einzelnen Durchlauf der Demo.
  \item \textbf{Metrik:} Der Diagonalfehler ist nicht direkt mit Winkelfehlern
        oder Zentimetern aus der Literatur vergleichbar.
\end{itemize}
